\documentclass[11pt, a4paper]{article}

% ── Fonts & encoding ─────────────────────────────────────────
\usepackage[T1]{fontenc}
\usepackage{lmodern}
\usepackage{microtype}
\usepackage{algorithm}
\usepackage{algpseudocode}

% ── Page geometry ────────────────────────────────────────────
\usepackage[
  top=2.5cm,
  bottom=2.5cm,
  left=2.5cm,
  right=2.5cm
]{geometry}

% ── Maths ────────────────────────────────────────────────────
\usepackage{amsmath}
\usepackage{amssymb}

% ── Tables & figures ─────────────────────────────────────────
\usepackage{booktabs}
\usepackage{float}
\usepackage{graphicx}
\usepackage{caption}
\captionsetup{
  font=small,
  labelfont=bf,
  skip=6pt
}

% ── Section formatting ───────────────────────────────────────
\usepackage{titlesec}
\titleformat{\section}
  {\large\bfseries}{\thesection.}{0.5em}{}
\titleformat{\subsection}
  {\normalsize\bfseries}{\thesubsection.}{0.5em}{}
\titlespacing*{\section}{0pt}{1.5ex plus .5ex minus .2ex}{1ex plus .2ex}
\titlespacing*{\subsection}{0pt}{1.2ex plus .4ex minus .2ex}{0.8ex plus .2ex}

% ── Headers & footers ───────────────────────────────────────
\usepackage{fancyhdr}
\pagestyle{fancy}
\fancyhf{}
\fancyhead[L]{\small\textit{ML: k-Means and PCA for Regime Detection}}
\fancyhead[R]{\small\textit{Jing et al.}}
\fancyfoot[C]{\thepage}
\renewcommand{\headrulewidth}{0.4pt}

% ── Hyperlinks ───────────────────────────────────────────────
\usepackage{hyperref}
\hypersetup{
  colorlinks=true,
  linkcolor=blue!70!black,
  citecolor=blue!70!black,
  urlcolor=blue!70!black
}

% ── Spacing ──────────────────────────────────────────────────
\usepackage{setspace}
\setstretch{1.15}
\setlength{\parskip}{0.4em}
\setlength{\parindent}{0pt}

% ═════════════════════════════════════════════════════════════
\title{\textbf{Machine Learning: k-Means and PCA\\[4pt] for Regime Detection}}
\author{Zhidian Jing, Siqi Zhang, Stefaan Fernando, Daniel Greenstein, Oscar Peyron\\[2pt]
\small University College London \quad -- \quad COMP0040}
\date{April 2026}

\begin{document}

\maketitle
\thispagestyle{fancy}

% ─────────────────────────────────────────────────────────────
\section{Introduction}

% ─────────────────────────────────────────────────────────────
\section{Data Processing and Feature Engineering}
This project focuses on regime detection on multi-dimensional financial time-series.
We track 9 major stock indices across 3 different regions (America, Europe, Asia) to develop our analysis.

\begin{table}[H]
\centering
\caption{Index universe by region.}
\label{tab:index_universe}
\renewcommand{\arraystretch}{1.3}
\begin{tabular}{l l l}
\toprule
\textbf{Region} & \textbf{Index} & \textbf{Ticker} \\
\midrule
America & S\&P 500    & \texttt{\^{}GSPC}  \\
America & Dow Jones   & \texttt{\^{}DJI}   \\
America & Nasdaq      & \texttt{\^{}IXIC}  \\
\addlinespace
Europe  & FTSE 100    & \texttt{\^{}FTSE}  \\
Europe  & CAC 40      & \texttt{\^{}FCHI}  \\
Europe  & DAX         & \texttt{\^{}GDAXI} \\
\addlinespace
Asia    & KOSPI       & \texttt{\^{}KS11}  \\
Asia    & Nikkei 225  & \texttt{\^{}N225}  \\
Asia    & Hang Seng   & \texttt{\^{}HSI}   \\
\bottomrule
\end{tabular}
\end{table}

\subsection{Raw Data Structure}
Each index is downloaded from Yahoo Finance with daily OHLCV columns:
\texttt{Date}, \texttt{Open}, \texttt{High}, \texttt{Low}, \texttt{Close}, and \texttt{Volume}.
After merging, the full DataFrame uses a MultiIndex on columns of the form
\texttt{(index\_name, OHLCV\_field)}, so that, for example, S\&P\,500 closing
prices are accessed via \texttt{df["SP500"]["Close"]}.

\subsection{Feature Engineering}
From the closing prices two families of features are computed.

\paragraph{Log-returns.}
Over multiple horizons (default: 1, 5, and 21 trading days):
\begin{equation}
\label{eq:log_return}
r_{t,p} = \ln\!\left(\frac{P_{t}}{P_{t-p}}\right).
\end{equation}

\paragraph{Rolling annualised volatility.}
Default windows of 21 and 63 days:
\begin{equation}
\label{eq:vol}
\sigma_{t,w} = \operatorname{std}\!\left(r_{t-w:t}\right) \times \sqrt{252}\,.
\end{equation}

This produces a feature matrix of shape
$\bigl(T,\; 9 \times (3 \text{ returns} + 2 \text{ volatilities})\bigr) = (T, 45)$
columns, where each row corresponds to one trading day and each column to a
pair \texttt{(index\_name, feature\_name)}.

In order to capture the dynamic nature of regime, features are computed in a rolling window fashion, so that the feature vector at time $t$ is based on data from the previous $w$ days.
For example, the rolling correlation between two indices $X$ and $Y$ at time $t$ with window $w$ is computed as:
\begin{equation}
\label{eq:rolling_corr}
    \rho_{X,Y,t} = \frac{\sum_{i=0}^{w-1} (r_{X,t-i} - \bar{r}_X)(r_{Y,t-i} - \bar{r}_Y)}{\sqrt{\sum_{i=0}^{w-1} (r_{X,t-i} - \bar{r}_X)^2 \;\sum_{i=0}^{w-1} (r_{Y,t-i} - \bar{r}_Y)^2}}
\end{equation}
where $\bar{r}_X$ and $\bar{r}_Y$ are the mean returns of indices $X$ and $Y$ over the window $w$.

We additionally compute momentum Z-score features to measure mean-reversion by normalising the distance between the current price and the moving average by the rolling volatility:
\begin{equation}
\label{eq:zscore}
    Z_{t,w} = \frac{P_t - \operatorname{MA}_{t,w}}{\sigma_{t,w}}
\end{equation}

\subsection{Pipeline Summary}

\begin{figure}[H]
\centering
\renewcommand{\arraystretch}{1.4}
\begin{tabular}{c}
\toprule
\textbf{Yahoo Finance API} \\
$\big\downarrow$ \\
Raw OHLCV (per index) \\
$\big\downarrow$ \\
Merged DataFrame (MultiIndex columns) \\
$\big\downarrow$ \\
Close prices only (simple DataFrame) \\
$\big\downarrow$ \\
\begin{tabular}{c@{\hspace{2em}}c}
Log-returns (1d, 5d, 21d) & Rolling volatility (21d, 63d) \\
\end{tabular} \\
$\big\downarrow$ \\
\textbf{Feature matrix $\rightarrow$ Clustering / PCA} \\
\bottomrule
\end{tabular}
\caption{Data pipeline from raw download to feature matrix.}
\label{fig:pipeline}
\end{figure}

% ─────────────────────────────────────────────────────────────
\section{Mathematical Foundations}
In this project we apply two main machine learning techniques for regime detection: k-means clustering and Principal Component Analysis (PCA).

\subsection{K-Means Clustering}
K-means clustering is an unsupervised learning algorithm that partitions the data into $K$ clusters by minimising the within-cluster sum of squares. The algorithm iteratively assigns each data point to the nearest cluster centroid and then updates the centroids based on the assigned points.

The Within-Cluster Sum of Squares (WCSS) is defined as:
\begin{equation}
\label{eq:wcss}
\text{WCSS} = \sum_{k=1}^{K} \sum_{x_i \in C_k} \| x_i - \mu_k \|^2
\end{equation}
where $C_k$ is the set of points in cluster $k$ and $\mu_k$ is the centroid of cluster $k$.

However, when dealing with higher-dimensional data, the standard k-means algorithm can be limited due to the curse of dimensionality as well as the different topologies of the clusters.
To address this, one can reduce dimensionality using PCA before applying k-means.

\subsection{Principal Component Analysis (PCA)}

% ─────────────────────────────────────────────────────────────
\section{Proposed Algorithm}
We propose the following algorithm for regime detection: 
\begin{algorithm}[H]
\caption{K-Means Regime Detection Pipeline}
\label{alg:regime_detection}
\begin{algorithmic}[1]

\Require Closing price matrix $\mathbf{P} \in \mathbb{R}^{T \times N}$ for $N$ indices over $T$ trading days
\Require PCA variance threshold $\alpha \in (0,1]$, candidate range $[K_{\min}, K_{\max}]$

\Ensure Regime labels $\ell_t \in \{0, \dots, K^{*}\!-\!1\}$ for each day $t$, ordered by mean return

\Statex
\Statex \textbf{— Stage 1: Feature Engineering —}
\For{each index $n = 1, \dots, N$}
    \For{each horizon $p \in \{1, 5, 21\}$}
        \State $r_{t,p}^{(n)} \gets \ln\!\bigl(P_t^{(n)} / P_{t-p}^{(n)}\bigr)$ \Comment{Log-returns}
    \EndFor
    \For{each window $w \in \{21, 63\}$}
        \State $\sigma_{t,w}^{(n)} \gets \operatorname{std}\!\bigl(r_{t\text{-}w:t}^{(n)}\bigr) \times \sqrt{252}$ \Comment{Rolling volatility}
    \EndFor
    \State $Z_t^{(n)} \gets \bigl(P_t^{(n)} - \operatorname{MA}_{t,63}^{(n)}\bigr) \big/ \sigma_{t,63}^{(n)}$ \Comment{Momentum z-score}
    \State $\text{RSI}_t^{(n)} \gets \text{RSI}_{14}\!\bigl(P^{(n)}\bigr)$ \Comment{Relative Strength Index}
\EndFor
\For{each pair $(i, j)$ with $i < j$}
    \State $\rho_{t}^{(i,j)} \gets \operatorname{RollingCorr}_{63}\!\bigl(r^{(i)}, r^{(j)}\bigr)$ \Comment{Pairwise correlation}
\EndFor
\State Concatenate all features into $\mathbf{F} \in \mathbb{R}^{T' \times D}$ and drop warm-up rows

\Statex
\Statex \textbf{— Stage 2: Standardisation —}
\For{each feature $d = 1, \dots, D$}
    \State $\tilde{F}_{t,d} \gets \bigl(F_{t,d} - \hat{\mu}_d\bigr) / \hat{\sigma}_d$ \Comment{Z-score normalisation}
\EndFor

\Statex
\Statex \textbf{— Stage 3: Dimensionality Reduction —}
\State $\mathbf{U}, \mathbf{\Sigma}, \mathbf{V}^\top \gets \operatorname{SVD}(\tilde{\mathbf{F}})$
\State Select $d^{*} = \min\bigl\{d : \sum_{i=1}^{d} \sigma_i^2 \big/ \sum_{i=1}^{D} \sigma_i^2 \geq \alpha \bigr\}$
\State $\mathbf{X} \gets \tilde{\mathbf{F}} \, \mathbf{V}_{:, 1:d^{*}} \in \mathbb{R}^{T' \times d^{*}}$ \Comment{Project onto top $d^{*}$ PCs}

\Statex
\Statex \textbf{— Stage 4: Optimal $K$ Selection —}
\For{$K = K_{\min}, \dots, K_{\max}$}
    \State Run K-Means on $\mathbf{X}$ with $K$ clusters
    \State Record inertia $\mathcal{I}(K)$ and silhouette score $\mathcal{S}(K)$
\EndFor
\State $K^{*} \gets \arg\max_{K} \; \mathcal{S}(K)$ \Comment{Best silhouette}

\Statex
\Statex \textbf{— Stage 5: Final Clustering —}
\State $\hat{\ell}_1, \dots, \hat{\ell}_{T'} \gets \operatorname{KMeans}(\mathbf{X},\; K^{*})$

\Statex
\Statex \textbf{— Stage 6: Regime Ordering —}
\For{each cluster $k = 0, \dots, K^{*}\!-\!1$}
    \State $\bar{r}_k \gets \operatorname{mean}\!\bigl(\{r_t : \hat{\ell}_t = k\}\bigr)$ \Comment{Mean basket return}
\EndFor
\State $\pi \gets \operatorname{argsort}\!\bigl(\bar{r}_0, \dots, \bar{r}_{K^{*}-1}\bigr)$ \Comment{Rank worst $\to$ best}
\State $\ell_t \gets \pi^{-1}(\hat{\ell}_t) \quad \forall\, t$ \Comment{Relabel: $0$ = bearish, $K^{*}\!-\!1$ = bullish}

\Statex
\State \Return $\boldsymbol{\ell} = (\ell_1, \dots, \ell_{T'})$

\end{algorithmic}
\end{algorithm}

% ─────────────────────────────────────────────────────────────
\section{Results}

% ─────────────────────────────────────────────────────────────
\section{Conclusion}

\end{document}